% sec-01: the week in one record.  Table 21, Figure 13 (graphviz-type record
% chain, one rank per weekday).
\section{The Week in One Record}

\subsection{Four signals, one a day}

Each signal of the week was answered on the day whose theme it fits: federal
program offices on Monday, regulators on Tuesday, the structures the company
works inside on Wednesday, and pilots and solicitations on Thursday. Table 21
sets the four beside the letters each produced and what each is waiting on. The
SEC office's words are the only words of any office quoted in this block, and
they are quoted exactly; no White House communication is quoted or paraphrased
anywhere \cite{whcontact2026}.

\begin{apptable}
\begin{tabularx}{\textwidth}{>{\raggedright\arraybackslash}p{1.9cm} >{\raggedright\arraybackslash}p{4.6cm} >{\raggedright\arraybackslash}p{3.6cm} >{\raggedright\arraybackslash}X}
\headrow \thead{Day} & \thead{Signal} & \thead{Sent} & \thead{Waiting on}\\
\midrule
Monday & Two unsolicited LinkedIn applicants & 3 emails, 2 replies & NIH, the SBA, UC San Diego; each applicant \\
Tuesday & The SEC office, ``very seriously'' & 5 emails & The SEC office and the FDA, on venue \\
Wednesday & White House communications & 4 emails, 1 form & OSTP, the Genesis Mission, Moores \\
Thursday & An NSF pilot inquiry & 4 emails, 2 filings & NSF on the Pitch; the pilot's gate \\
Friday & None & 3 follow-ups & The same answers, each with a date \\
\end{tabularx}
\end{apptable}
\vspace{-0.60cm}
\tabcap{Table 21. The week's four signals, one a day, each with what it produced and what it\\
is waiting on; Friday spends no signal and sends only the follow-ups that are earned.}

\subsection{The same week, drawn as a chain}

Figure 13 draws the week as a chain of records, one rank per weekday, from the
signal to the letters it produced to the answers still pending. The bold edges
on the right are the follow-ups: Monday's and Tuesday's pending answers reach
Friday's letters through the rule, and Thursday's letter to NIH joins Monday's
in one consolidated follow-up. Wednesday's and Thursday's other letters reach
nothing on Friday, because their interval has not yet run: theirs fall due on
Monday and Tuesday of the following week.

\begin{appfloat}[!tb]
\begin{appfig}[graphviz-type / record chain, one rank per weekday]
% Five ranks left to right.  In each: a header cell, the signal as an ellipse,
% the letters as a two-field record, and the pending answers as an ellipse.
% The follow-up edges run in a channel beneath the ranks and rise into
% Friday's record from the left.
\foreach \x/\d in {0.00/Monday, 3.20/Tuesday, 6.40/Wednesday, 9.60/Thursday, 12.80/Friday}{
  \node[gvcellh,text width=22mm,minimum height=5mm] at (\x,0.05) {\d};}
\node[gvkey,text width=19mm] (s1) at (0.00,-0.95) {Two unsolicited applicants};
\node[gvkey,text width=19mm] (s2) at (3.20,-0.95) {The SEC office's update};
\node[gvkey,text width=19mm] (s3) at (6.40,-0.95) {White House communications};
\node[gvkey,text width=19mm] (s4) at (9.60,-0.95) {An NSF pilot inquiry};
\node[gvgray,text width=19mm] (s5) at (12.80,-0.95) {No new signal};
\foreach \i/\x/\a/\b in {1/0.00/{3 emails, 2 replies}/{Personnel, first hire, training},
  2/3.20/{5 emails}/{Venue; 3 notices ask nothing},
  3/6.40/{4 emails, 1 form}/{A contributor's place},
  4/9.60/{4 emails, 2 filings}/{Gate, resources, disclosure}}{
  \node[gvcells,text width=25mm,minimum height=5mm] (l\i a) at (\x,-2.05) {\a};
  \node[gvcell,text width=25mm,minimum height=5mm] (l\i b) at (\x,-2.57) {\b};}
\node[gvcellh,text width=25mm,minimum height=5mm] (l5a) at (12.80,-2.05) {3 follow-ups};
\node[gvcell,text width=25mm,minimum height=5mm] (l5b) at (12.80,-2.57) {Only what has earned one};
\node[gvgray,text width=19mm] (p1) at (0.00,-3.80) {NIH, SBA, UC San Diego; each applicant};
\node[gvgray,text width=19mm] (p2) at (3.20,-3.80) {The SEC office and the FDA, on venue};
\node[gvgray,text width=19mm] (p3) at (6.40,-3.80) {OSTP, Genesis, Moores; due Monday};
\node[gvgray,text width=19mm] (p4) at (9.60,-3.80) {NSF, NCTI, NAIRR; due Tuesday};
\node[gvsoft,text width=19mm] (p5) at (12.80,-3.80) {Read first on Monday at 5:00};
\foreach \i in {1,2,3,4,5}{
  \draw[gvedge] (s\i) -- (l\i a);
  \draw[gvedge] (l\i b) -- (p\i);}
\draw[gvedgeb] (l4a.east) -- node[above=1pt,font=\tiny,inner sep=1pt] {joins} (l5a.west);
\draw[gvundirb] (p1.south) -- (0.00,-4.85) -- (11.20,-4.85);
\draw[gvundirb] (p2.south) -- (3.20,-4.85);
\draw[gvedgeb] (11.20,-4.85) -- (11.20,-2.57) -- (l5b.west);
\node[font=\tiny,fill=fundwhite,inner sep=1pt] at (6.40,-4.85) {follow-up, on the third open business day};
\ptitle{-1.40}{0.85}{From signal to letter to pending answer}
\end{appfig}
\vspace{-0.60cm}
\figcaption{Figure 13. The week as a record chain, one rank per weekday, with the follow-up edges\\
that let Monday's and Tuesday's letters, and no others, reach Friday's three follow-ups.}
\end{appfloat}

\subsection{The totals}

Sixteen emails went out on the first four days, each to at least three verified
addresses, drawn from a register of thirty-six. Two LinkedIn replies went on
Monday, and two more are drafted today, each to be sent only if its applicant has
answered. Three submissions were made: the White House contact form, the NSF
Project Pitch, and the Research.gov organization record \cite{nsf26510}. Two
words are attributed to the SEC office, and none to any other office.

\subsection{The standing fact}

The trial behind every letter of the week is built around daraxonrasib, which
the FDA approved on August 26, 2026 as Rasonque for metastatic pancreatic
adenocarcinoma after prior systemic therapy
\cite{fdarasonque2026,fdadisco2026,revmedapproval2026}. The approval followed
Revolution Medicines' Phase 1/2 program and its Phase 3 RASolute 302 trial
\cite{ctgov001,ctgov302,rasolute302}. ChemicalQDevice's AI papers on the
molecule were developed independently of that program and simultaneously with
it, from the June 2025 identification study \cite{daraxident} to the Phase 1 and
Phase 2 protocols and the investigational new drug application of 2026
\cite{onpremwhippl,phase2proto,phase1ind}. The perioperative use the company's
protocol proposes remains investigational.
